\documentclass[aspectratio=169,10pt]{beamer}

\usepackage{iftex}
\providecommand{\dunenoirsetupfonts}{}
\newcommand{\DuneNoirLocalFontPath}{fonts}

\newcommand{\DuneNoirMonoFallback}{%
  \IfFontExistsTF{JetBrainsMono Nerd Font Propo}{
    \setmonofont{JetBrainsMono Nerd Font Propo}[Scale=0.92]
  }{
    \IfFontExistsTF{JetBrainsMono Nerd Font}{
      \setmonofont{JetBrainsMono Nerd Font}[Scale=0.92]
    }{
      \setmonofont{Courier New}[Scale=0.92]
    }
  }
}

\newcommand{\DuneNoirSansFallback}{%
  \IfFontExistsTF{IBM Plex Sans}{
    \setsansfont{IBM Plex Sans}
    \setmainfont{IBM Plex Sans}
  }{
    \IfFontExistsTF{Helvetica Neue}{
      \setsansfont{Helvetica Neue}
      \setmainfont{Helvetica Neue}
    }{
      \setsansfont{Arial}
      \setmainfont{Arial}
    }
  }
}

\newcommand{\DuneNoirUseHelvetica}{%
  \renewcommand{\dunenoirsetupfonts}{%
    \IfFontExistsTF{Helvetica Neue}{
      \setsansfont{Helvetica Neue}
      \setmainfont{Helvetica Neue}
    }{
      \IfFontExistsTF{Helvetica}{
        \setsansfont{Helvetica}
        \setmainfont{Helvetica}
      }{
        \DuneNoirSansFallback
      }
    }
    \DuneNoirMonoFallback
  }%
}

\newcommand{\DuneNoirUseOpenSans}{%
  \renewcommand{\dunenoirsetupfonts}{%
    \IfFileExists{\DuneNoirLocalFontPath/OpenSans.ttf}{
      \setsansfont{Open Sans Custom}[
        Path=\DuneNoirLocalFontPath/,
        UprightFont={OpenSans.ttf},
        ItalicFont={OpenSans-Italic.ttf},
        BoldFont={OpenSans.ttf},
        BoldItalicFont={OpenSans-Italic.ttf},
        BoldFeatures={FakeBold=1.45},
        BoldItalicFeatures={FakeBold=1.45}
      ]
      \setmainfont{Open Sans Custom}[
        Path=\DuneNoirLocalFontPath/,
        UprightFont={OpenSans.ttf},
        ItalicFont={OpenSans-Italic.ttf},
        BoldFont={OpenSans.ttf},
        BoldItalicFont={OpenSans-Italic.ttf},
        BoldFeatures={FakeBold=1.45},
        BoldItalicFeatures={FakeBold=1.45}
      ]
    }{
      \DuneNoirSansFallback
    }
    \DuneNoirMonoFallback
  }%
}

\newcommand{\DuneNoirUseDosis}{%
  \renewcommand{\dunenoirsetupfonts}{%
    \IfFileExists{\DuneNoirLocalFontPath/Dosis.ttf}{
      \setsansfont{Dosis Custom}[
        Path=\DuneNoirLocalFontPath/,
        UprightFont={Dosis.ttf},
        ItalicFont={Dosis.ttf},
        BoldFont={Dosis.ttf},
        BoldItalicFont={Dosis.ttf},
        ItalicFeatures={FakeSlant=0.16},
        BoldFeatures={FakeBold=1.6},
        BoldItalicFeatures={FakeBold=1.6,FakeSlant=0.16}
      ]
      \setmainfont{Dosis Custom}[
        Path=\DuneNoirLocalFontPath/,
        UprightFont={Dosis.ttf},
        ItalicFont={Dosis.ttf},
        BoldFont={Dosis.ttf},
        BoldItalicFont={Dosis.ttf},
        ItalicFeatures={FakeSlant=0.16},
        BoldFeatures={FakeBold=1.6},
        BoldItalicFeatures={FakeBold=1.6,FakeSlant=0.16}
      ]
    }{
      \DuneNoirSansFallback
    }
    \DuneNoirMonoFallback
  }%
}

\newcommand{\DuneNoirUseKoHo}{%
  \renewcommand{\dunenoirsetupfonts}{%
    \IfFileExists{\DuneNoirLocalFontPath/KoHo-Regular.ttf}{
      \setsansfont{KoHo Custom}[
        Path=\DuneNoirLocalFontPath/,
        UprightFont={KoHo-Regular.ttf},
        BoldFont={KoHo-Bold.ttf},
        ItalicFont={KoHo-Italic.ttf},
        BoldItalicFont={KoHo-BoldItalic.ttf}
      ]
      \setmainfont{KoHo Custom}[
        Path=\DuneNoirLocalFontPath/,
        UprightFont={KoHo-Regular.ttf},
        BoldFont={KoHo-Bold.ttf},
        ItalicFont={KoHo-Italic.ttf},
        BoldItalicFont={KoHo-BoldItalic.ttf}
      ]
    }{
      \DuneNoirSansFallback
    }
    \DuneNoirMonoFallback
  }%
}

\DuneNoirUseHelvetica
\makeatletter
\input{../dune_dark_template_lab/beamerthemeDUNENoir.sty}
\makeatother

\ifPDFTeX
  \usepackage[T1]{fontenc}
  \usepackage[utf8]{inputenc}
\fi

\usepackage{amsmath,amssymb}
\usepackage{booktabs,array}
\usepackage{graphicx}
\usepackage{hyperref}
\usepackage{listings}
\usepackage{tikz}
\usetikzlibrary{arrows.meta,calc,positioning,fit,shapes.geometric}

\title[DAPHNE DAQ Integration]{DAPHNE DAQ Integration}
\subtitle{Interface contract, DAQ data path, calibration scans, and mitigation evidence}
\author{Manuel Arroyave}
\date{June 2026}
\newcommand{\mydate}{June 2026}

\lstdefinestyle{terminal}{
  basicstyle=\ttfamily\scriptsize,
  columns=fullflexible,
  keepspaces=true,
  breaklines=true,
  frame=single,
  rulecolor=\color{noirline},
  backgroundcolor=\color{noirpanel},
  xleftmargin=0.2em,
  xrightmargin=0.2em
}

\newcommand{\sectioncard}[2]{%
  \begin{frame}[plain]
    \begin{tikzpicture}[remember picture,overlay]
      \fill[noirpanelalt] (current page.south west) rectangle (current page.north east);
      \fill[duneorange] ($(current page.north west)+(0.9,-0.9)$) rectangle ++(1.8,0.08);
      \fill[noircyan] ($(current page.south east)+(-2.8,0.9)$) rectangle ++(1.8,0.08);
      \draw[noirline, line width=0.8pt] ($(current page.south west)+(0.9,1.1)$) -- ($(current page.south east)+(-0.9,1.1)$);
    \end{tikzpicture}
    \vspace*{0.21\paperheight}
    {\color{noirmuted}\small\bfseries SECTION\par}
    \vspace{3mm}
    {\color{white}\bfseries\fontsize{29}{32}\selectfont #1\par}
    \vspace{4mm}
    {\color{noirmuted}\fontsize{13}{16}\selectfont #2\par}
  \end{frame}
}

\newcommand{\eyebrow}[1]{{\color{noircyan}\small\bfseries #1\par}}

\newcommand{\panelbox}[3][panel]{%
  \begin{beamercolorbox}[wd=\linewidth,sep=6pt,rounded=true]{#1}
    {\color{noircyan}\bfseries #2\par}
    \vspace{0.5mm}
    {\small #3\par}
  \end{beamercolorbox}
}

\newcommand{\statcard}[2]{%
  \begin{beamercolorbox}[wd=\linewidth,sep=6pt,rounded=true]{panel alt}
    {\color{noircyan}\bfseries\fontsize{14}{16}\selectfont #1\par}
    \vspace{0.6mm}
    {\color{noirmuted}\scriptsize #2\par}
  \end{beamercolorbox}
}

\newcommand{\blocksep}{%
  \vspace{0.7mm}%
  {\color{noirline}\hrule height 0.35pt}%
  \vspace{0.7mm}%
}

\newcommand{\resourcebar}[4]{%
  \begin{tikzpicture}[x=0.1cm,y=0.34cm,baseline=(label.base)]
    \pgfmathsetmacro{\barend}{#2/10}
    \node[anchor=east,font=\footnotesize,text=noirmuted] (label) at (0,0) {#1};
    \draw[draw=noirline,fill=noirpanelalt,rounded corners=1pt] (0.25,-0.17) rectangle (10.25,0.17);
    \draw[draw=#4,fill=#4!30,rounded corners=1pt] (0.25,-0.17) rectangle ({0.25+\barend},0.17);
    \node[anchor=west,font=\footnotesize,text=noirtext] at (10.45,0) {#3};
  \end{tikzpicture}\par
}

\tikzset{
  flowbox/.style={
    draw=noirline,
    fill=noirpanelalt,
    text=noirtext,
    rounded corners=4pt,
    align=center,
    text width=2.7cm,
    minimum height=0.92cm,
    inner sep=4pt,
    font=\scriptsize
  },
  flowboxwide/.style={
    draw=noirline,
    fill=noirpanelalt,
    text=noirtext,
    rounded corners=4pt,
    align=center,
    text width=3.35cm,
    minimum height=0.92cm,
    inner sep=4pt,
    font=\scriptsize
  },
  repo/.style={
    draw=noircyan!70,
    fill=noirpanel,
    text=noirtext,
    rounded corners=5pt,
    align=left,
    text width=3.55cm,
    minimum height=2.1cm,
    inner sep=7pt,
    font=\scriptsize
  },
  arrow/.style={-Latex, line width=0.8pt, draw=noircyan},
  warnarrow/.style={-Latex, line width=0.8pt, draw=duneorange},
  smallnote/.style={align=center, font=\tiny, text=noirtext}
}

\newcommand{\daqPipelineDiagram}{%
\begin{tikzpicture}[node distance=0.38cm and 0.32cm,
  flowboxwide/.append style={font=\footnotesize, text width=2.35cm, minimum height=1.02cm},
  smallnote/.append style={font=\scriptsize}]
  \node[flowboxwide] (hw) {DAPHNE\\Ethernet output};
  \node[flowboxwide, right=of hw] (daq) {DUNE DAQ\\readout path};
  \node[flowboxwide, right=of daq] (frag) {Fragments\\with DAPHNE ETH frames};
  \node[flowboxwide, right=of frag] (decode) {Python/C++\\decode interface};
  \node[flowboxwide, right=of decode] (analysis) {HDF5/NumPy\\Waffles analysis};
  \draw[arrow] (hw) -- (daq);
  \draw[arrow] (daq) -- (frag);
  \draw[arrow] (frag) -- (decode);
  \draw[arrow] (decode) -- (analysis);
  \node[smallnote, below=0.25cm of daq] {data taking baseline};
  \node[smallnote, below=0.25cm of decode] {decode contract};
\end{tikzpicture}
}

\newcommand{\standaloneDiagram}{%
\resizebox{0.9\linewidth}{!}{%
\begin{tikzpicture}[node distance=0.55cm and 0.55cm]
  \node[flowbox] (scripts) {\texttt{daphne-server}\\scan scripts};
  \node[flowbox, right=of scripts] (tunnel) {SSH tunnel\\or direct host};
  \node[flowbox, right=of tunnel] (server) {DAPHNE control\\server};
  \node[flowbox, right=of server] (spy) {Spy buffer\\or osc capture};
  \node[flowbox, below=0.85cm of server] (ssp) {SSP LED\\control};
  \node[flowbox, right=of spy] (products) {run folders\\plots, JSON};
  \draw[arrow] (scripts) -- (tunnel);
  \draw[arrow] (tunnel) -- (server);
  \draw[arrow] (server) -- (spy);
  \draw[arrow] (spy) -- (products);
  \draw[warnarrow] (scripts) |- (ssp);
  \draw[warnarrow] (ssp) -| (spy);
  \node[smallnote, above=0.22cm of server] {no DAQ run control};
  \node[smallnote, below=0.22cm of products] {calibration evidence};
\end{tikzpicture}
}%
}

\newcommand{\pdsRunDiagram}{%
\resizebox{0.94\linewidth}{!}{%
\begin{tikzpicture}[node distance=0.48cm and 0.42cm]
  \node[flowbox] (paths) {\texttt{00\_paths}\\facility};
  \node[flowbox, right=of paths] (defaults) {\texttt{02}/\texttt{03}\\defaults};
  \node[flowbox, right=of defaults] (scan) {scan JSON\\one variable};
  \node[flowbox, right=of scan] (patch) {minimal DAPHNE\\JSON patch};
  \node[flowbox, right=of patch] (xml) {\texttt{add\_daphne\_conf}\\OKS XML};
  \node[flowbox, right=of xml] (ssp) {\texttt{set\_ssp\_conf}\\LED};
  \node[flowbox, right=of ssp] (drunc) {\texttt{drunc}\\short run};
  \draw[arrow] (paths) -- (defaults);
  \draw[arrow] (defaults) -- (scan);
  \draw[arrow] (scan) -- (patch);
  \draw[arrow] (patch) -- (xml);
  \draw[arrow] (xml) -- (ssp);
  \draw[arrow] (ssp) -- (drunc);
  \node[smallnote, below=0.2cm of xml] {overlay, not hand edit};
  \node[smallnote, below=0.2cm of drunc] {data + monitoring};
\end{tikzpicture}
}%
}

\begin{document}

{
\setbeamertemplate{headline}{}
\setbeamertemplate{footline}{}
\begin{frame}[plain,noframenumbering]
  \begin{tikzpicture}[remember picture,overlay]
    \fill[noirbg] (current page.south west) rectangle (current page.north east);
    \fill[noirpanelalt] (current page.north west) rectangle ($(current page.north east)+(0,-2.1)$);
    \fill[duneorange] ($(current page.north west)+(0.8,-0.62)$) rectangle ++(1.5,0.08);
    \fill[noircyan] ($(current page.south east)+(-2.8,0.8)$) rectangle ++(1.9,0.08);
    \IfFileExists{logos/DUNElogo_white.png}{%
      \node[anchor=south east,opacity=0.13] at ($(current page.south east)+(-0.55,1.48)$) {%
        \includegraphics[width=0.33\paperwidth]{logos/DUNElogo_white.png}%
      };
    }{}%
    \IfFileExists{logos/FNAL-Logo-NAL-Light.png}{%
      \node[anchor=south east] at ($(current page.south east)+(-0.92,0.82)$) {%
        \includegraphics[height=0.45cm]{logos/FNAL-Logo-NAL-Light.png}%
      };
    }{}%
  \end{tikzpicture}
  \vspace*{0.26\paperheight}
  {\color{white}\bfseries\fontsize{29}{31}\selectfont DAPHNE DAQ Integration\par}
  \vspace{2mm}
  {\color{noirmuted}\fontsize{15}{18}\selectfont Interface contract, DAQ data path,\\ calibration scans, and mitigation evidence\par}
  \vspace{7mm}
  {\color{noircyan}\bfseries Manuel Arroyave\par}
  \vspace{1mm}
  {\color{noirmuted}June 2026\par}
\end{frame}
}

\begin{frame}[t]{Executive Claim}
\begin{columns}[T,onlytextwidth]
\column{0.49\textwidth}
\eyebrow{INTEGRATION STATUS}
{\color{white}\bfseries\fontsize{14}{17}\selectfont
DAPHNE Ethernet data taking is established in DAQ. Selected calibration scans are moving into \texttt{pds-run}. \par}

\vspace{1.5mm}
{\color{noirmuted}\small
EDMS 2088726 v7.0 anchors the DAQ/PDS interface; DAQ libraries own decode, with Waffles downstream. \par}

\column{0.46\textwidth}
    \vspace{-1.8cm}
\begin{beamercolorbox}[wd=\linewidth,sep=7pt,rounded=true]{panel alt}
\small
{\color{noircyan}\bfseries Interface}\par
    EDMS 2088726 v7.0 is the interface reference.\par \vspace{2mm}
\blocksep \vspace{2mm}

{\color{noircyan}\bfseries Nov 2025}\par
DAPHNE Ethernet data taking proven in the VD Coldbox.\par \vspace{2mm}
\blocksep \vspace{2mm}

{\color{noircyan}\bfseries May 2026}\par
Selected DAPHNE settings exercised through DAQ runs.\par \vspace{2mm}
\blocksep \vspace{2mm}

{\color{noircyan}\bfseries Firmware evidence}\par
Peak descriptors and grouped-Hermes studies support mitigation.\par
\end{beamercolorbox}
\end{columns}
\end{frame}

\begin{frame}[t]{Integration Evidence}
\begin{columns}[T,onlytextwidth]
\column{0.31\textwidth}
\panelbox{Formal interface}{%
Reference the DAQ/PDS boundary through EDMS 2088726 v7.0, for both VD and HD detectors.}

\column{0.31\textwidth}
\panelbox{DAQ data path}{%
DAPHNE Ethernet frames enter DAQ, are decoded through DAQ libraries, and are consumed downstream.}

\column{0.31\textwidth}
\panelbox{DAQ-wrapped scans}{%
Selected DAPHNE settings are scanned by \texttt{pds-run}: every scan point becomes a DAQ run. Naturally, future calibrations become datasets within a run.}
\end{columns}

\vspace{2mm}
\begin{beamercolorbox}[wd=\linewidth,sep=6pt,rounded=true]{panel alt}
\footnotesize
\textcolor{noircyan}{\textbf{Integration scope:}} interface reference, DAQ formats and decode, controller/monitoring path, and run-controlled DAPHNE scans. Resource and dead-time mitigation are summarized on dedicated evidence slides.
\end{beamercolorbox}
\end{frame}

\begin{frame}[t]{Interface Contract}
\begin{columns}[T,onlytextwidth]
\column{0.42\textwidth}
\eyebrow{FORMAL REFERENCE}
{\color{white}\bfseries DUNE Interface Control Document: DAQ \& Photon Detector System (PDS)\par}
\vspace{2mm}
{\color{noircyan}\bfseries\fontsize{18}{21}\selectfont
EDMS 2088726 v7.0\\
\par}
\vspace{1mm}
{\color{noirmuted}\footnotesize
\href{https://edms.cern.ch/document/2088726/7}{https://edms.cern.ch/document/2088726/7}\par}

\vspace{2mm}
\begin{beamercolorbox}[wd=\linewidth,sep=7pt,rounded=true]{panel alt}
\small
This is the common DAQ/PDS boundary for the readout, timing, control, monitoring, and data-format discussion.
\end{beamercolorbox}

\column{0.54\textwidth}
\begin{beamercolorbox}[wd=\linewidth,sep=8pt,rounded=true]{panel}
\small
\textcolor{noircyan}{\textbf{Interface scope}}\par
\vspace{1mm}
\begin{itemize}
  \item DAQ/PDS boundary and responsibilities.
  \item DAPHNE-to-DAQ 10 GbE UDP data interface.
  \item DAQ header and DAPHNE header responsibilities.
  \item Timing provenance and calibration-command tagging.
  \item CCM configuration, control, monitoring, and error reporting path.
\end{itemize}
\end{beamercolorbox}

\end{columns}
\end{frame}

\sectioncard{DAQ Baseline}{The implementation maps the ICD data/control boundary into DAQ packages, decode utilities, and monitoring}

\begin{frame}[t]{Baseline Data Path}
\centering
\vspace{2mm}
\daqPipelineDiagram

\vspace{3mm}
\noindent
\begin{minipage}[t]{0.30\textwidth}
\panelbox{Format}{C++ classes describe the wire layout and enforce bit-field boundaries.}
\end{minipage}\hfill
\begin{minipage}[t]{0.30\textwidth}
\panelbox{Decode}{Python bindings expose arrays for analysis without rewriting the unpacker.}
\end{minipage}\hfill
\begin{minipage}[t]{0.30\textwidth}
\panelbox{Waffles}{Waffles consumes the decoded arrays for calibration and validation workflows.}
\end{minipage}
\end{frame}

\begin{frame}[t]{Changed DAQ Repositories}
\vspace{1mm}
\noindent
\begin{minipage}[t]{0.31\textwidth}
\begin{beamercolorbox}[wd=\linewidth,sep=7pt,rounded=true]{panel}
\small
{\color{noircyan}\bfseries fddetdataformats\par}
\vspace{1mm}
Frame definitions and unit tests for DAPHNE Ethernet data.\par
\vspace{1mm}
\textcolor{duneorange}{Peak descriptor support included.}
\end{beamercolorbox}
\end{minipage}\hfill
\begin{minipage}[t]{0.31\textwidth}
\begin{beamercolorbox}[wd=\linewidth,sep=7pt,rounded=true]{panel alt}
\small
{\color{noircyan}\bfseries rawdatautils\par}
\vspace{1mm}
DAQ-side unpackers expose ADC, timestamp, channel, stream, and peak arrays.\par
\vspace{1mm}
\textcolor{noirmuted}{Used by validation and Waffles.}
\end{beamercolorbox}
\end{minipage}\hfill
\begin{minipage}[t]{0.31\textwidth}
\begin{beamercolorbox}[wd=\linewidth,sep=7pt,rounded=true]{panel}
\small
{\color{noircyan}\bfseries daphnemodules\par}
\vspace{1mm}
DAPHNE V3 controller path for configuration and operational monitoring.\par
\vspace{1mm}
\textcolor{noirmuted}{Includes V3 \texttt{GeneralInfo} opmon.}
\end{beamercolorbox}
\end{minipage}

\vspace{3mm}
\begin{beamercolorbox}[wd=\linewidth,sep=6pt,rounded=true]{panel alt}
\footnotesize
\textcolor{noircyan}{\textbf{Together:}} format, decode, and controller changes are separate DAQ packages, but they form one DAPHNE readout chain.
\end{beamercolorbox}
\end{frame}

\begin{frame}[t]{Data-Format Contract}
\begin{columns}[T,onlytextwidth]
\column{0.5\textwidth}
\begin{beamercolorbox}[wd=\linewidth,sep=3.5pt,rounded=true]{panel}
\footnotesize
{\color{noircyan}\bfseries\texttt{DAPHNEEthFrame}\par}
Single-channel frame: Ethernet header, DAPHNE metadata, five peak descriptors, and 1024 packed 14-bit ADC samples.
\end{beamercolorbox}

\vspace{0.8mm}
\begin{beamercolorbox}[wd=\linewidth,sep=3.5pt,rounded=true]{panel alt}
\footnotesize
{\color{noircyan}\bfseries Peak descriptors\par}
New firmware information is now persisted in the frame and exposed through DAQ formats and Python decode utilities.
\end{beamercolorbox}

\vspace{0.8mm}
\begin{beamercolorbox}[wd=\linewidth,sep=3.5pt,rounded=true]{panel}
\footnotesize
{\color{noircyan}\bfseries\texttt{DAPHNEEthStreamFrame}\par}
Four-channel stream frame with 280 samples per channel packed as 14-bit ADCs.
\end{beamercolorbox}

\column{0.46\textwidth}
\begin{beamercolorbox}[wd=\linewidth,sep=3.5pt,rounded=true]{panel alt}
\footnotesize
\textcolor{noircyan}{\textbf{Contract checks}}
\vspace{0.3mm}
\begin{itemize}
  \item ICD-backed constants: ADC width, sample count, channel count.
  \item Accessors validate channel, sample, and peak indices.
  \item Unit tests round-trip ADC and peak fields.
  \item Bindings register triggered and stream types.
\end{itemize}
\end{beamercolorbox}
\end{columns}
\end{frame}

\begin{frame}[t]{Decode Contract}
\begin{columns}[T,onlytextwidth]
\column{0.49\textwidth}
\eyebrow{RAW DATA UTILS}
\small
\texttt{rawdatautils} is not a second format definition. It reuses \texttt{fddetdataformats} and exposes arrays consumed by Waffles and validation scripts:

\begin{itemize}
  \item ADC, timestamp, and channel arrays
  \item stream-frame equivalents
  \item peak descriptor arrays for triggered frames
\end{itemize}

\column{0.5\textwidth}
\begin{tikzpicture}[node distance=0.3cm]
  \node[flowboxwide] (frag) {DAQ fragment};
  \node[flowboxwide, below=of frag] (reinterpret) {reinterpret as\\DAPHNE frame};
  \node[flowboxwide, below=of reinterpret] (arrays) {NumPy arrays\\ADC, time, channel};
  \node[flowboxwide, below=of arrays] (analysis) {Waffles\\analysis and validation};
  \draw[arrow] (frag) -- (reinterpret);
  \draw[arrow] (reinterpret) -- (arrays);
  \draw[arrow] (arrays) -- (analysis);
\end{tikzpicture}
\end{columns}
\end{frame}

\begin{frame}[t]{Controller And Monitoring Contract}
\begin{columns}[T,onlytextwidth]
\column{0.5\textwidth}
\panelbox{\texttt{daphnemodules} V3 controller}{%
Configures the DAPHNE V3 analog chain, reads counters, and publishes opmon in DAQ.}

\vspace{1mm}
\panelbox[panel alt]{Important release finding}{%
\texttt{fddaq-v5.5.0-a9-1} ships 3.0.1. V3 \texttt{GeneralInfo} opmon appears in 3.0.2 on \texttt{patch/fddaq-v5.5.x}.}

\column{0.46\textwidth}
\begin{tikzpicture}[node distance=0.3cm]
  \node[flowboxwide] (counter) {ReadTriggerCounters};
  \node[flowboxwide, below=of counter] (perchan) {per-channel\\threshold, record, busy, full};
  \node[flowboxwide, below=of perchan] (general) {ReadGeneralInfo\\bias, power, temp};
  \node[flowboxwide, below=of general] (opmon) {Kafka opmon\\shifter view};
  \draw[arrow] (counter) -- (perchan);
  \draw[arrow] (perchan) -- (general);
  \draw[arrow] (general) -- (opmon);
\end{tikzpicture}
\end{columns}
\end{frame}

\begin{frame}[t]{Dead-Time Mitigation Evidence}
\vspace{-3mm}
\begin{columns}[T,onlytextwidth]
\column{0.21\textwidth}
\eyebrow{OPERATING POINT}
{\color{white}\bfseries\fontsize{13}{15}\selectfont 4.7 kHz/channel\par}
{\color{noirmuted}\tiny rounded comparison}

\column{0.20\textwidth}
\begin{beamercolorbox}[wd=\linewidth,sep=3pt,rounded=true]{panel}
{\color{noircyan}\bfseries Main\par}
{\color{white}\bfseries\fontsize{18}{20}\selectfont 7\%\par}
{\color{noirmuted}\tiny 1024-tick waveform}
\end{beamercolorbox}

\column{0.20\textwidth}
\begin{beamercolorbox}[wd=\linewidth,sep=3pt,rounded=true]{panel alt}
{\color{noircyan}\bfseries Mitigation\par}
{\color{white}\bfseries\fontsize{18}{20}\selectfont 2.5\%\par}
{\color{noirmuted}\tiny 512-tick waveform}
\end{beamercolorbox}

\column{0.33\textwidth}
\begin{beamercolorbox}[wd=\linewidth,sep=3pt,rounded=true]{panel}
\scriptsize
\textcolor{noircyan}{\textbf{Main caveat}}\par
The lower point assumes waveform records are restricted to 512 ticks.
\end{beamercolorbox}
\end{columns}

\vspace{0.8mm}
\centering
\begin{tikzpicture}[
  topobox/.style={
    draw=noirline,
    fill=noirpanelalt,
    text=noirtext,
    rounded corners=4pt,
    align=center,
    text width=1.9cm,
    minimum height=0.66cm,
    inner sep=3.5pt,
    font=\scriptsize
  },
  emphbox/.style={topobox, draw=duneorange, fill=noirpanel},
  newbox/.style={topobox, draw=noircyan, fill=noirpanel},
  topoarrow/.style={-Latex, line width=0.7pt, draw=noircyan}
]
  \node[anchor=west,font=\bfseries\scriptsize,text=noircyan] at (-6.05,0.67) {Main HDL topology};
  \node[topobox] (mch) at (-5.0,-0.08) {40 channel\\samples};
  \node[topobox, right=0.35cm of mch] (mst) {40 self-trigger\\paths};
  \node[emphbox, right=0.35cm of mst] (mprim) {trigger\\primitives\\+ builders};
  \node[emphbox, right=0.35cm of mprim] (mmux) {2 readout\\muxes};
  \node[topobox, right=0.35cm of mmux] (mhermes) {1 Hermes\\10 Gb link};
  \draw[topoarrow] (mch) -- (mst);
  \draw[topoarrow] (mst) -- (mprim);
  \draw[topoarrow] (mprim) -- (mmux);
  \draw[topoarrow] (mmux) -- (mhermes);

  \node[anchor=west,font=\bfseries\scriptsize,text=noircyan] at (-6.05,-0.9) {Mitigation HDL topology};
  \node[topobox] (nch) at (-5.0,-1.65) {40 channel\\samples};
  \node[newbox, right=0.35cm of nch] (ngroup) {4 channels\\per source};
  \node[newbox, right=0.35cm of ngroup] (nsrc) {10 grouped\\sources};
  \node[newbox, right=0.35cm of nsrc] (ninputs) {10 Hermes\\inputs};
  \node[topobox, right=0.35cm of ninputs] (nlink) {1 Hermes\\10 Gb link};
  \draw[topoarrow] (nch) -- (ngroup);
  \draw[topoarrow] (ngroup) -- (nsrc);
  \draw[topoarrow] (nsrc) -- (ninputs);
  \draw[topoarrow] (ninputs) -- (nlink);
\end{tikzpicture}
\end{frame}

\begin{frame}[t]{K26 Resource Mitigation Evidence}
\begin{columns}[T,onlytextwidth]
\column{0.49\textwidth}
\eyebrow{BASELINE, OUTPUT-2C337B5}
\vspace{1.5mm}
\resourcebar{CLB sites}{99.55}{99.55\%}{duneorange}
\resourcebar{DSP48E2}{99.36}{99.36\%}{duneorange}
\resourcebar{CLB LUTs}{74.48}{74.48\%}{noircyan}
\resourcebar{Block RAM}{68.75}{68.75\%}{noircyan}
\resourcebar{URAM}{62.50}{62.50\%}{noircyan}
\resourcebar{Registers}{44.78}{44.78\%}{noircyan}

\column{0.49\textwidth}
\eyebrow{GROUPED CANDIDATE, E5E2D70}
\vspace{1.5mm}
\resourcebar{CLB LUTs}{76.05}{76.05\%}{noircyan}
\resourcebar{Block RAM}{67.36}{67.36\%}{noircyan}
\resourcebar{RAMB36/FIFO}{63.89}{63.89\%}{noircyan}
\resourcebar{URAM}{62.50}{62.50\%}{noircyan}
\resourcebar{Registers}{61.19}{61.19\%}{noircyan}
\resourcebar{DSP48E2}{0.00}{0.00\%}{noircyan}
\end{columns}

\vspace{2mm}
\begin{beamercolorbox}[wd=\linewidth,sep=6pt,rounded=true]{panel alt}
\footnotesize
The grouped candidate keeps LUT, BRAM, and URAM below device limits and removes the baseline DSP pressure. The remaining qualification is readout behavior, not FPGA fit or route legality.
\end{beamercolorbox}
\end{frame}

\sectioncard{Calibration Path Into DAQ}{Selected scans already generate DAQ configurations, DAQ runs, and reviewable datasets}

\begin{frame}[t]{Calibration Scans As DAQ Artifacts}
\vspace{-1mm}
\eyebrow{CURRENT STATUS}
{\color{white}\bfseries\fontsize{13.5}{16}\selectfont
Selected calibration variables now generate DAQ configurations and run-associated datasets. \par}

\vspace{2.5mm}
\begin{columns}[T,onlytextwidth]
\column{0.31\textwidth}
\begin{beamercolorbox}[wd=\linewidth,sep=5pt,rounded=true]{panel}
\footnotesize
\textcolor{noircyan}{\textbf{Calibration probes}}\par
\vspace{0.4mm}
Probing identified variables, ranges, and target settings worth moving into DAQ.
\end{beamercolorbox}

\column{0.31\textwidth}
\begin{beamercolorbox}[wd=\linewidth,sep=5pt,rounded=true]{panel alt}
\footnotesize
\textcolor{noircyan}{\textbf{DAQ wrapper}}\par
\vspace{0.4mm}
\texttt{pds-run} drafts the DAQ workflow: generate the overlay and delegate each scan point to \texttt{drunc}.
\end{beamercolorbox}

\column{0.31\textwidth}
\begin{beamercolorbox}[wd=\linewidth,sep=5pt,rounded=true]{panel}
\footnotesize
\textcolor{noircyan}{\textbf{Dataset path}}\par
\vspace{0.4mm}
The result is DAQ data and monitoring tied to the configuration that produced it.
\end{beamercolorbox}
\end{columns}

\vspace{2mm}
{\color{noirmuted}\scriptsize
Integration point: the mechanics are functional now; the endpoint is the same scan semantics implemented natively inside DAQ.\par}
\end{frame}

\begin{frame}[t]{Draft DAQ Workflow For Variable Scans}
\centering
\vspace{2mm}
\pdsRunDiagram

\vspace{3mm}
\noindent
\begin{minipage}[t]{0.23\textwidth}
\panelbox{Scan plan}{One variable, explicit values, and optional SSP LED settings.}
\end{minipage}\hfill
\begin{minipage}[t]{0.23\textwidth}
\panelbox{DAQ config}{Minimal DAPHNE overlay generated for the scan point.}
\end{minipage}\hfill
\begin{minipage}[t]{0.23\textwidth}
\panelbox{\texttt{drunc} wrapper}{Each scan point delegates to DAQ run control.}
\end{minipage}\hfill
\begin{minipage}[t]{0.23\textwidth}
\panelbox{Dataset}{Data and monitoring are tied to the DAQ run.}
\end{minipage}
\end{frame}

\begin{frame}[t]{Functional Scan Coverage}
\vspace{2mm}
\noindent
\begin{minipage}[t]{0.23\textwidth}
\panelbox{Self-trigger}{Threshold values become short DAQ runs with changed trigger settings.}
\end{minipage}\hfill
\begin{minipage}[t]{0.23\textwidth}
\panelbox{AFE response}{Attenuation, bias, offset, and trim become generated overlays plus run data.}
\end{minipage}\hfill
\begin{minipage}[t]{0.23\textwidth}
\panelbox{LED / SSP}{LED intensity and mask are coordinated with DAQ data taking.}
\end{minipage}\hfill
\begin{minipage}[t]{0.23\textwidth}
\panelbox{Plan mode}{Generated XML and run-control actions can be inspected before execution.}
\end{minipage}

\vspace{3mm}
{\color{noirmuted}\small
The important result is not command syntax: variable scans already exercise the DAQ configuration path and produce run-associated datasets with the current system.\par}
\end{frame}

\begin{frame}[t]{Operational Environment}
\begin{columns}[T,onlytextwidth]
\column{0.48\textwidth}
\panelbox{Login path}{%
\texttt{marroyav@np04-srv-017} then \texttt{np04daq@np04-srv-024}; use tmux session \texttt{np02-vdcb}.}
\vspace{1mm}
\panelbox{Workarea}{%
\texttt{/nfs/sw/dunedaq/}\\
\texttt{dunedaq-fddaq-v5.5.0-dev-pds}.\\
Source \texttt{setup\_for\_run.sh}.}

\column{0.48\textwidth}
\panelbox{Proxy state}{%
Use \texttt{source \textasciitilde/bin/web\_proxy.sh} only for package downloads. Disable for running with \texttt{-u}.}
\vspace{1mm}
\panelbox{Hardware target}{%
Board identity and IP must be checked in the active JSON/XML; target DAPHNE selection is a configuration requirement, not an assumption.}
\end{columns}
\end{frame}

\sectioncard{Validation And Review}{The success condition is both software cleanliness and live DAQ behavior}

\begin{frame}[t]{Live Validation Result}
\begin{columns}[T,onlytextwidth]
\column{0.24\textwidth}
\statcard{Run 44355}{Default release without the \texttt{daphnemodules} override: data taking worked, but V3 monitoring was not visible.}
\column{0.24\textwidth}
\statcard{Run 44356}{Clean \texttt{daphnemodules} 3.0.2 override: two-minute TEST run succeeded at 20 Hz.}
\column{0.24\textwidth}
\statcard{Firmware}{Peak descriptors were added in firmware and then carried into the DAQ data-format contract.}
\column{0.24\textwidth}
\statcard{Conclusion}{Monitoring loss was traced to the release package version, not to run control.}
\end{columns}

\vspace{3mm}
\begin{beamercolorbox}[wd=\linewidth,sep=6pt,rounded=true]{panel alt}
\footnotesize
Final tested state: \texttt{sourcecode/daphnemodules} commit \texttt{13289b5} on \texttt{patch/fddaq-v5.5.x}; \texttt{DAPHNE\_V3\_ROUTE} unset; monitoring visible.
\end{beamercolorbox}
\end{frame}

\begin{frame}[t]{Live \texttt{pds-run} Scan Exercise}
\vspace{-1mm}
\begin{columns}[T,onlytextwidth]
\column{0.58\textwidth}
\eyebrow{SCAN MODES EXERCISED}
\vspace{0.8mm}
\begin{beamercolorbox}[wd=\linewidth,sep=5pt,rounded=true]{panel}
\footnotesize
\renewcommand{\arraystretch}{1.22}
\begin{tabular}{@{}>{\raggedright\arraybackslash}p{0.32\linewidth}>{\raggedright\arraybackslash}p{0.57\linewidth}@{}}
\textbf{Self-trigger} & thresholds: 11, 10 \\
\textbf{AFE attenuation} & all attenuators: 2200, 2300 \\
\textbf{Offset} & selected channels: 2260, 2275 \\
\textbf{Trim} & selected channels: 1, 0 \\
\textbf{AFE bias} & AFE0: 10, 0; others fixed at 0 \\
\textbf{LED / SSP} & intensity: 3000, 4095; mask 4 \\
\end{tabular}
\end{beamercolorbox}

\vspace{1.1mm}
{\color{noirmuted}\footnotesize Run span: 44357--44368; each point generated a DAQ configuration and delegated execution to \texttt{drunc}.\par}

\column{0.38\textwidth}
\eyebrow{RESULT}
\vspace{0.8mm}
\begin{beamercolorbox}[wd=\linewidth,sep=5pt,rounded=true]{panel alt}
\small
Successful scan commands exited with status 0.
\blocksep
\texttt{xmllint --noout} passed after each XML update.
\blocksep
No successful-run traceback or \texttt{drunc} command failure was found.
\end{beamercolorbox}

\vspace{0.8mm}
\begin{beamercolorbox}[wd=\linewidth,sep=4pt,rounded=true]{panel}
\footnotesize
\textcolor{noircyan}{\textbf{Evidence location}}\par
\texttt{np04-srv-024}; branch \texttt{vdcb-daq-scans}.\\
Logs: \texttt{/tmp/pds\_vdcb\_live\_*}.
\end{beamercolorbox}
\end{columns}
\end{frame}

\begin{frame}[t]{Live Exercise Findings}
\begin{beamercolorbox}[wd=\linewidth,sep=4pt,rounded=true]{panel}
\scriptsize
\renewcommand{\arraystretch}{1.16}
\begin{tabular}{@{}>{\raggedright\arraybackslash}p{0.25\linewidth}>{\raggedright\arraybackslash}p{0.65\linewidth}@{}}
\toprule
\textbf{Finding} & \textbf{Consequence} \\
\midrule
SSP object name & Active XML uses \texttt{np02-ssp}; scan configs must not silently rely on the old \texttt{np02-ssp-on} default. \\
Execution seed & \texttt{99\_apply\_fe\_min\_v2\_subset.json} is plan-safe but not execution-safe; \texttt{add\_daphne\_conf} requires \texttt{detector\_id}. \\
Zero/default XML & \texttt{trim=0}, \texttt{v\_bias=0}, and \texttt{bias\_ctrl=0} can be omitted by \texttt{add\_daphne\_conf}; absence of these attributes is not by itself a run failure. \\
Post-scan state & The exercise left SSP at \texttt{channel\_mask=4}, \texttt{pulse1\_width=3}; scan plans should state whether to restore the baseline. \\
\bottomrule
\end{tabular}
\end{beamercolorbox}

\vspace{1.5mm}
\begin{beamercolorbox}[wd=\linewidth,sep=4pt,rounded=true]{panel alt}
\scriptsize
Conclusion: the DAQ-wrapped scan loop is operational for the tested modes; remaining cleanup is configuration hygiene, explicit target naming, and post-scan restoration policy.
\end{beamercolorbox}
\end{frame}

\begin{frame}[t]{Summary}
\vspace{1mm}
\begin{columns}[T,onlytextwidth]
\column{0.52\textwidth}
\eyebrow{ESTABLISHED FOR PRR}
\vspace{0.8mm}
\begin{beamercolorbox}[wd=\linewidth,sep=6pt,rounded=true]{panel}
\small
\textbf{Interface:} EDMS 2088726 v7.0 anchors the DAQ/PDS boundary for VD and HD.
\blocksep
\textbf{DAQ path:} DAPHNE Ethernet data taking, format ownership, and decode utilities have been exercised.
\blocksep
\textbf{Calibration path:} selected DAPHNE settings can be scanned through \texttt{pds-run}; each point becomes a DAQ run.
\end{beamercolorbox}

\column{0.43\textwidth}
\eyebrow{BOUNDED QUALIFICATION}
\vspace{0.8mm}
\begin{beamercolorbox}[wd=\linewidth,sep=6pt,rounded=true]{panel alt}
\small
\textbf{Dead time:} mitigation evidence gives 7\% to 2.5\% at 4.7 kHz/channel, with 512-tick records.
\blocksep
\textbf{Resources:} grouped Hermes strategy addresses the K26 fit pressure.
\blocksep
\textbf{Next check:} Hermes readout, backpressure, and real-trigger stress.
\end{beamercolorbox}
\end{columns}
\end{frame}

\end{document}
